% year: תשע"ט
% semester: א
% moed: ב
% number: 3ב
% points: 10
% categories: derivatives, proofs, true-false
% source: exams/מבחנים/תשעט/מועד ב תשעט.pdf

\begin{question}
(10 נק') הוכיחו או הפריכו: אם פונקציה $f(x)$ גזירה בנקודה $x_0$ היא רציפה ב-$x_0$.
\end{question}

\begin{hint}
כתבו $f(x)-f(x_0)=\frac{f(x)-f(x_0)}{x-x_0}\cdot(x-x_0)$.
\end{hint}

\begin{solution}
\textbf{הטענה נכונה.} נניח ש-$f$ גזירה ב-$x_0$, כלומר הגבול $f'(x_0)=\lim_{x\to x_0}\frac{f(x)-f(x_0)}{x-x_0}$ קיים וסופי. לכל $x\ne x_0$ בסביבת $x_0$:
\[ f(x)-f(x_0)=\frac{f(x)-f(x_0)}{x-x_0}\cdot(x-x_0). \]
לפי אריתמטיקה של גבולות (מכפלת שני גבולות סופיים):
\[ \lim_{x\to x_0}\big(f(x)-f(x_0)\big)=f'(x_0)\cdot0=0, \]
כלומר $\lim_{x\to x_0}f(x)=f(x_0)$, ו-$f$ רציפה ב-$x_0$. $\blacksquare$

(הכיוון ההפוך אינו נכון: $|x|$ רציפה ב-$0$ אך אינה גזירה שם.)
\end{solution}
